\documentclass[10pt]{article}

\usepackage[margin=0.85in]{geometry}
\usepackage{array}
\usepackage{longtable}
\usepackage{hyperref}

\title{QuorumScope: A Bounded Consistency Testing Workbench for Replicated Registers}
\author{QuorumScope artifact report}
\date{\today}

\newcommand{\qs}{QuorumScope}
\newcommand{\firstack}{first-ack}
\newcommand{\quorum}{quorum}
\newcommand{\cmd}[1]{\texttt{#1}}

\begin{document}
\maketitle

\begin{abstract}
\qs{} is a bounded consistency-testing workbench for a replicated single-key register under
network partitions. The artifact compares a deliberately unsafe \firstack{} protocol, which
accepts the first reachable replica response, against a majority \quorum{} protocol that may
return unavailable instead of stale data. The core engine includes a deterministic partition
simulator, concurrent operation histories, a linearizability checker, counterexample shrinking,
seeded adversarial search, a manifest-backed regression corpus, and a tiny bounded exhaustive
explorer. The current evaluation finds stale-read violations for \firstack{} in curated,
generated, and exhaustive corpora. Under the same modeled assumptions, \quorum{} produced zero
linearizability violations in the declared bounded corpora, while reporting unavailable
operations as an explicit availability tradeoff. This report is an artifact technical report:
it documents the implemented model, reproducible commands, observed results, and limitations. It
does not claim a formal proof for arbitrary replicated systems.
\end{abstract}

\section{Introduction}

Replicated storage systems can acknowledge an operation locally while still violating a global
consistency property. A small partition can make this failure mode concrete: a majority side
accepts a write, while a minority side later returns an older value. The local read may look
successful because some replicas responded, but the resulting history is not linearizable if the
write completed before the read began.

\qs{} is built to make that failure inspectable. Rather than only displaying a hand-written
split-brain example, the artifact can generate partition schedules, execute competing protocols
on the same schedule, check the resulting histories, shrink failing scenarios, preserve selected
counterexamples as regression fixtures, and enumerate a tiny finite model exhaustively. The
workbench also exposes the checker state so that a user can see why a read is accepted, blocked
by real-time order, or rejected by the single-register oracle.

The goal is bounded engineering evidence, not a broad production model. The useful claim is not
``quorum is proven correct.'' The useful claim is narrower: for the declared fixtures, default
adversarial generated corpus, and tiny exhaustive model reported below, no \quorum{}
linearizability violations were observed under the current register and partition assumptions,
and the availability cost was counted.

\section{System Model}

The model is a deterministic replicated register with a single key and a finite set of replicas.
Clients issue read and write operations. Each successful operation records an invocation time,
a response time, the reachable replicas consulted by the protocol, the returned value or written
value, and a status. Operations that cannot complete under a protocol are recorded as
unavailable. Unavailable operations are part of the product result because they expose the
safety/availability tradeoff, but they are not treated as successful completed operations in the
linearizability check.

Partitions are represented as reachability groups over replica identifiers. A scenario is a
bounded sequence of operation and topology steps. The simulator uses deterministic timing, so
the same scenario, seed, and configuration reproduce the same event trace. The current model
does not simulate a real network stack, arbitrary message delay interleavings, durability,
leader election, retries, or message loss beyond the modeled reachability constraints.

The artifact includes two protocols:

\begin{itemize}
  \item \textbf{\firstack{}.} A read or write completes after the first reachable replica response.
  This is intentionally unsafe under partitions because it can read from a minority side that has
  not observed a completed write.
  \item \textbf{\quorum{}.} A read or write requires a majority of replicas. In the modeled
  scenarios, this prevents the stale minority read from completing, but it may make operations
  unavailable while a partition is active.
\end{itemize}

Both protocols are run against the same generated or enumerated schedules so that the comparison
is about protocol behavior rather than different workloads.

\section{Linearizability Checking}

The checker consumes a history of completed operations over the single-key register. It enforces
real-time order: if operation $a$ finishes before operation $b$ starts, any legal sequential
linearization must place $a$ before $b$. It then searches for a sequential order that preserves
those constraints and satisfies the register specification. Writes update the oracle value; reads
must return the current oracle value at their linearization point.

If a legal order exists, the checker returns a witness order and final register value. If no legal
order exists, the checker returns a violation witness when available. A common witness is a stale
read: a read returned value $v_0$ even though a completed write to $v_1$ must precede it in every
legal order. For example, the default curated split-brain replay reports:
\begin{quote}
\texttt{op3 read returned v0 after op2 write completed with v1.}
\end{quote}

The checker also emits oracle diagnostics. These diagnostics include the successful operations
checked, unavailable operations ignored by the linearizability search, real-time predecessor
constraints, the number of search states explored, memoized dead-end hits, candidate operations
per captured state, and candidate statuses. Candidate statuses distinguish operations that are
ready, blocked by real-time predecessors, or rejected because a read does not match the current
oracle value.

The checker proves only a property of the supplied finite history: whether that history is
linearizable with respect to the implemented single-register specification. It does not prove a
protocol correct for all histories, all networks, or all implementations.

\section{Counterexample Search and Shrinking}

\qs{} has three paths for producing histories.

First, the curated split-brain replay demonstrates a minimal human-readable stale read: a write
completes on the majority side, then a minority read under \firstack{} returns the old value.

Second, seeded adversarial search generates bounded partition schedules with configurable seed,
operation count, client count, read/write mix, partition intensity, and concurrency intensity.
The default search run used in this report explored 50 seeds from seed 143 with 5 nodes, 8
operations, 3 clients, partition intensity 0.75, and concurrent intensity 0.45. It found a first
\firstack{} violation at seed 143, with an original 11-step scenario shrunk to 3 steps:
\begin{quote}
\texttt{op2 read returned v0 after op1 write completed with v143-0-x1.}
\end{quote}
The reproduction command printed by the tool is:
\begin{verbatim}
npm run search -- --seed 143 --seeds 1 --protocol compare --nodes 5 \
  --ops 8 --clients 3 --read-ratio 0.55 --chaos 0.75 \
  --concurrency 0.45 --shrink
\end{verbatim}

Third, selected minimized failures are promoted into a manifest-backed regression corpus. The
current corpus contains four fixtures: a curated split-brain stale read, a minimized adversarial
search failure, a concurrent safe-overlap fixture, and a minimized exhaustive stale-read witness.
The corpus runner validates expectations and reports mismatches if fixture behavior drifts.

\section{Bounded Exhaustive Explorer}

The bounded exhaustive explorer enumerates a deliberately tiny finite model. The current default
configuration uses 3 replicas, 2 clients, a single key, at most 3 operations, at most 2 topology
changes, healed topology plus canonical 1/2 partitions, and deterministic simulator timing.
Concurrent schedules are included, but message timing interleavings are not exhaustively
enumerated.

For each enumerated terminal history, both protocols are checked. A terminal history is one
fully generated scenario at the configured operation and topology bounds, after simulator
execution and history extraction. The explorer reports exactly what it covered: prefixes
explored, terminal histories checked, unique scenarios, concurrent schedules, partition shapes,
operation patterns, violations, stale-read witnesses, and unavailable operations.

In the current run, the explorer checked 1000 terminal histories and explored 1339 prefixes.
It found 144 \firstack{} violations, including 126 stale-read witnesses. For \quorum{}, it found
0 linearizability violations and counted 1064 unavailable operations. The bounded claim emitted
by the tool is intentionally scoped:
\begin{quote}
Within the tiny 3-replica / 2-client / 3-operation scenario model explored here, quorum produced
0 linearizability violations under the current register and partition assumptions. This is not a
proof for arbitrary systems.
\end{quote}

\section{Evaluation}

All numbers in this section come from commands run on the current checkout before writing this
report. The product verification command completed successfully and includes tests, typecheck,
production build, UI smoke, demo, corpus replay, search comparison, exhaustive exploration, and
product report.

\begin{table}[h]
\centering
\begin{tabular}{|p{0.31\linewidth}|p{0.6\linewidth}|}
\hline
\textbf{Command} & \textbf{Observed result} \\
\hline
\cmd{npm test} & 16 test files passed; 82 tests passed. \\
\hline
\cmd{npm run verify:product} & Product verification passed; UI smoke reported 20/20 checks. \\
\hline
\cmd{npm run corpus} & 4 fixtures; 8 expected outcomes matched; 0 mismatches. \\
\hline
\cmd{npm run search} & 50 seeds explored; first failing seed 143; 50 \firstack{} violations; 0 \quorum{} violations; 155 \quorum{} unavailable operations. \\
\hline
\cmd{npm run exhaustive} & 1000 terminal histories checked; 1339 prefixes explored; 144 \firstack{} violations; 0 \quorum{} violations; 1064 \quorum{} unavailable operations. \\
\hline
\cmd{npm run bench} & Deterministic benchmark seed 4310, 50 runs: \firstack{} had 50 violations and 50 stale-read witnesses; \quorum{} had 0 violations and 50 unavailable operations. \\
\hline
\end{tabular}
\caption{Current artifact verification snapshot.}
\end{table}

\begin{table}[h]
\centering
\small
\begin{tabular}{|p{0.25\linewidth}|r|r|r|p{0.23\linewidth}|}
\hline
\textbf{Surface} & \textbf{FA viol.} & \textbf{Q viol.} & \textbf{Q unavail.} & \textbf{Notes} \\
\hline
Regression corpus & 3/4 & 0/4 & 4 & 4 fixtures, 8 outcomes matched. \\
\hline
Adversarial search & 50/50 & 0/50 & 155 & First failure seed 143; 3 minimized steps. \\
\hline
Tiny exhaustive model & 144 & 0 & 1064 & 1000 terminal histories; 392 concurrent schedules. \\
\hline
Deterministic benchmark & 50/50 & 0/50 & 50 & Seed 4310; average 42 events. \\
\hline
\end{tabular}
\caption{Safety and availability results by product surface.}
\end{table}

The result pattern is consistent across surfaces: \firstack{} is unsafe under the modeled
partitions, while \quorum{} avoids the stale-read violations in these bounded corpora by refusing
some operations.

\section{Workbench and Product Surface}

The browser workbench is implemented with Vite, React, and TypeScript. It is a technical
workbench rather than a marketing page. The UI exposes the generated or loaded scenario,
protocol selection, replay controls, partition visualization, operation timeline, event trace,
checker verdict, minimized counterexample, protocol comparison, benchmark table, and
linearizability oracle diagnostics. The oracle panel shows real-time predecessor constraints and
captured checker search states so that the hard part of the artifact is visible.

The CLI is the primary reproducibility surface. \cmd{npm run report} aggregates corpus replay,
default adversarial search, tiny exhaustive exploration, protocol comparison, provenance, bounded
claims, and reproduction commands. \cmd{npm run smoke:ui} is a headless build smoke that checks
the generated Vite shell and core workbench surfaces; it is useful for unattended automation but
is not a substitute for interactive browser testing.

\section{Limitations}

\qs{} is intentionally narrow. The current artifact has the following limitations:

\begin{itemize}
  \item It models a single-key register, not a multi-key database or transaction system.
  \item The network model is a simplified partition/reachability model, not a real network.
  \item Search is bounded and generated; exhaustive exploration covers only the declared tiny
  finite model.
  \item Simulator timing is deterministic and does not enumerate all message delay schedules.
  \item The protocols are simple \firstack{} and majority \quorum{} register protocols, not Raft,
  Paxos, leader election, or a production replication system.
  \item The model does not include durability, storage crashes, retries, read repair, Byzantine
  faults, or real deployment behavior.
  \item The absence of \quorum{} violations in the reported corpora is bounded evidence under the
  current assumptions, not a general proof of quorum correctness.
\end{itemize}

\section{Future Work}

The most useful extensions would deepen the model without hiding its assumptions. Candidate
directions include a richer scheduler and message-delay model, stronger exhaustive bounds for
small configurations, multi-key histories, predicate reads or transaction-isolation variants,
better interactive replay and import/export workflows, and additional protocol models once the
current register model remains stable under the expanded checks. Any future protocol additions
should include explicit assumptions, deterministic fixtures, corpus expectations, and bounded
claims.

\appendix
\section{Reproducibility Commands}

The following commands reproduce the current artifact surfaces from a clean checkout with Node
dependencies installed:

\begin{verbatim}
npm install
npm test
npm run verify:product
npm run report
npm run bench
npm run search
npm run exhaustive
npm run corpus
npm run smoke:ui
\end{verbatim}

The key counterexample reproduction commands emitted by the current tools are:
\begin{verbatim}
npm run search -- --seed 143 --seeds 1 --protocol compare --nodes 5 \
  --ops 8 --clients 3 --read-ratio 0.55 --chaos 0.75 \
  --concurrency 0.45 --shrink

npm run exhaustive -- --case ex-000043 --max-ops 3 --topology 2 \
  --clients 2 --seed 7001 --show
\end{verbatim}

\end{document}
